% ============================================================================
% 컴파일 방법 (Compile with XeLaTeX, NOT pdflatex):
%   xelatex textbook_GMM_DOB_Stability.tex
%   xelatex textbook_GMM_DOB_Stability.tex   (두 번 실행하여 목차 갱신)
%
% 본문은 chapters/ 폴더의 개별 .tex 파일을 \input 으로 포함한다.
% 원본은 Part_*.md (루트 폴더) — pandoc 으로 변환된다:
%   pandoc -f markdown -t latex --top-level-division=chapter \
%          Part_X_ChN_*.md -o chapters/Part_X_ChN_*.tex
% ============================================================================
\documentclass[11pt,a4paper,openright]{book}

% --- 패키지 ---------------------------------------------------------------
\usepackage{fontspec}
\usepackage{amsmath,amssymb,amsthm}
\usepackage{listings}
\usepackage{xcolor}
\usepackage{booktabs}
\usepackage{longtable}
\usepackage{array}
\usepackage{calc}
\usepackage{graphicx}
\usepackage{hyperref}
\usepackage{geometry}
\usepackage{tcolorbox}
\usepackage{enumitem}
\usepackage{etoolbox}

% --- 폰트 (fonts/ 하위의 Nanum TTF 사용) --------------------------------
\setmainfont{NanumMyeongjo}[
  Path = fonts/,
  Extension = .ttf,
  UprightFont = *-Regular,
  BoldFont = *-Bold ]
\setsansfont{NanumGothic}[
  Path = fonts/,
  Extension = .ttf,
  UprightFont = *-Regular,
  BoldFont = *-Bold ]
\setmonofont{NanumGothic}[
  Path = fonts/,
  Extension = .ttf,
  UprightFont = *-Regular,
  BoldFont = *-Bold,
  Scale = 0.92 ]

% 한글 줄바꿈
\XeTeXlinebreaklocale "ko"
\XeTeXlinebreakskip = 0pt plus 1pt

% --- 여백 -----------------------------------------------------------------
\geometry{
  inner=2.5cm, outer=2.5cm, top=2.5cm, bottom=2.5cm,
  headheight=1cm, headsep=0.5cm, footskip=0.5cm
}

% --- Theorems -------------------------------------------------------------
\theoremstyle{definition}
\newtheorem{definition}{정의}[chapter]
\newtheorem{theorem}[definition]{정리}
\newtheorem{lemma}[definition]{보조정리}
\newtheorem{corollary}[definition]{따름정리}
\newtheorem{proposition}[definition]{명제}
\newtheorem{remark}[definition]{주목}
\newtheorem{example}[definition]{예제}

% --- 코드 리스팅 ----------------------------------------------------------
\lstset{
  basicstyle=\ttfamily\small,
  keywordstyle=\color{blue},
  commentstyle=\color{gray},
  stringstyle=\color{red},
  breaklines=true,
  breakatwhitespace=true,
  showspaces=false,
  numbers=left,
  numberstyle=\tiny,
  stepnumber=1,
  tabsize=2,
  xleftmargin=10pt,
  extendedchars=true,
  inputencoding=utf8,
  literate={~}{{\textasciitilde}}1
}

% --- Pandoc 호환 매크로 ---------------------------------------------------
% pandoc이 생성하는 \tightlist, \passthrough, \real, quote 환경 등을 정의
\providecommand{\tightlist}{%
  \setlength{\itemsep}{0pt}\setlength{\parskip}{0pt}}
\providecommand{\passthrough}[1]{#1}
\providecommand{\real}[1]{#1}
% pandoc이 쓰는 'quote' 환경을 blockquote tcolorbox 로 대체
\renewenvironment{quote}
  {\begin{tcolorbox}[colback=blue!5,colframe=blue!75,arc=0.1cm,
     left=6pt,right=6pt,top=4pt,bottom=4pt]\itshape}
  {\end{tcolorbox}}

% hyperref 경고 완화: 유니코드 북마크
\hypersetup{
  unicode=true,
  colorlinks=true,
  linkcolor=blue!60!black,
  urlcolor=blue!50!black,
  citecolor=blue!50!black,
  pdftitle={GMM-DOB Stability Textbook},
  pdfauthor={DGIST MCL},
  bookmarksnumbered=true
}

% --- 제목 페이지 ----------------------------------------------------------
\title{\LARGE GMM 기반 로봇 제어기의 안정성 분석\\[4pt]
  \Large LMI에서 Fuzzy Lyapunov까지}
\author{대구경북과학기술원(DGIST) 모션제어연구실(MCL)}
\date{2024--2026}

\begin{document}

\frontmatter
\maketitle

\begin{center}
\vspace*{2cm}
\begin{minipage}{0.85\textwidth}
\begin{center}
\large
\textbf{GMM 기반 로봇 제어기의 안정성 분석}\\[2pt]
\normalsize LMI에서 Fuzzy Lyapunov까지
\end{center}
\vspace{1cm}
\small
이 교재는 Gaussian Mixture Model(GMM)로부터 학습된 로봇 제어기의
\emph{안정성}을 수학적으로 증명하는 전 과정을 다룬다.
Lyapunov 함수로부터 시작하여 LMI, SDP, Takagi--Sugeno Fuzzy 시스템,
Fuzzy Lyapunov Function, Descriptor Form, 그리고 DOB 보상항이 포함된
완전한 UUB(Uniformly Ultimately Bounded) 증명에 이르기까지
\textbf{수식적 엄밀성}과 \textbf{문헌 참조}를 동시에 만족하도록 작성되었다.
\end{minipage}
\end{center}

\tableofcontents

\mainmatter

% ============================================================================
% PART I — 기초 도구
% ============================================================================
\part{기초 도구: Lyapunov, LMI, SDP}

\input{chapters/Part_I_Ch1_Lyapunov}
\input{chapters/Part_I_Ch2_LMI}
\input{chapters/Part_I_Ch3_SDP_Solver}

% ============================================================================
% PART II — Fuzzy 안정성 이론
% ============================================================================
\part{Fuzzy 안정성 이론}

\input{chapters/Part_II_Ch4_TS_Fuzzy}
\input{chapters/Part_II_Ch5_FLF}
\input{chapters/Part_II_Ch6_UUB}

% ============================================================================
% PART III — GMM-DOB 안정성의 완전한 증명
% ============================================================================
\part{GMM-DOB 안정성의 완전한 증명}

\input{chapters/Part_III_Ch7_GMR_TS}
\input{chapters/Part_III_Ch8_Descriptor}
\input{chapters/Part_III_Ch9_Full_Proof}
\input{chapters/Part_III_Ch10_Design}

\backmatter

\end{document}
